
\section{TRATAMIENTO DEL CAMPO ACÚSTICO Y EXTRACCIÓN DE OBSERVABLES FÍSICOS}

En este capítulo se describe el tratamiento de las señales acústicas desde su adquisición en dos posiciones del espacio hasta la obtención de observables empleados en los modelos de localización y estimación cinemática. La exposición sigue la cadena de medición: propagación del campo, transducción, muestreo, acondicionamiento digital y extracción de magnitudes temporales y espectrales.

El propósito no es introducir una colección de operaciones digitales aisladas, sino establecer qué información física proporciona cada etapa, bajo qué hipótesis puede interpretarse y cuáles son sus limitaciones. En particular, se distingue entre la resolución nominal de los parámetros instrumentales y la incertidumbre efectiva de las magnitudes estimadas.

\paragraph{Nota sobre implementación y reproducibilidad.}
La implementación digital asociada al instrumento WALO se encuentra documentada en dos repositorios de acceso abierto:
\begin{itemize}
    \item \textbf{Lógica digital de adquisición y procesamiento en FPGA:}\\
    \url{https://github.com/SiririComun/hw-xadc-dma-overlays}
    \cite{SiririHW2026}.

    \item \textbf{Controladores y rutinas de procesamiento:}\\
    \url{https://github.com/SiririComun/sw-pynq-sound-localizer}
    \cite{SiririSW2026}.
\end{itemize}
Los repositorios permiten consultar la implementación de hardware y software. La reproducibilidad de un resultado experimental concreto requiere, además, disponer de la configuración utilizada, los datos originales y los procedimientos de análisis correspondientes.

% ==============================================================================
% 1.1. DEFINICIÓN DEL CAMPO ACÚSTICO
% ==============================================================================

\subsection{Definición del campo acústico}

Se considera la propagación de perturbaciones acústicas en el aire, tratado como un medio homogéneo e isótropo en primera aproximación. En ausencia de fuentes dentro de la región de observación y para perturbaciones de pequeña amplitud, la fluctuación de presión sonora $p(\vec{r},t)$ satisface la ecuación escalar de onda \cite{kinsler2000fundamentals, morse1968theoretical, pierce1989acoustics}:

\begin{equation}
    \nabla^2 p(\vec{r},t)
    -
    \frac{1}{c^2}
    \frac{\partial^2 p(\vec{r},t)}{\partial t^2}
    = 0,
    \label{eq:onda_acustica}
\end{equation}

donde $c$ es la velocidad de propagación del sonido en el medio. Para aire seco, puede emplearse como aproximación en el intervalo usual de temperaturas de laboratorio la expresión:

\begin{equation}
    c(T) =
    331.3
    \sqrt{1+\frac{T}{273.15}}
    \quad \si{\metre\per\second},
    \label{eq:velocidad_sonido_temperatura}
\end{equation}

con $T$ expresada en grados Celsius. A una temperatura de referencia de $\SI{20}{\celsius}$ se obtiene:

\begin{equation}
    c(20\,^\circ\mathrm{C})
    \approx \SI{343.2}{\metre\per\second}.
\end{equation}

Este valor es nominal: la velocidad del sonido depende de las condiciones del medio y, para una determinación experimental precisa de distancias o velocidades, debe utilizarse la temperatura medida durante el ensayo. La humedad y la composición del aire también pueden introducir variaciones respecto de la aproximación anterior.

El campo se observa mediante dos transductores ubicados en posiciones fijas $\vec{r}_1$ y $\vec{r}_2$. La separación entre ambos define la línea de base instrumental:

\begin{equation}
    d = \left\|\vec{r}_2-\vec{r}_1\right\|.
    \label{eq:linea_base}
\end{equation}

La magnitud de $d$ depende de la configuración experimental. En consecuencia, se mantiene como parámetro geométrico general y se especifica en los modelos y experimentos donde interviene.

Las señales acústicas en las posiciones de los receptores se expresan como:

\begin{equation}
    p_1(t)=p(\vec{r}_1,t),
    \qquad
    p_2(t)=p(\vec{r}_2,t).
    \label{eq:presion_receptores}
\end{equation}

Estas señales contienen la contribución de la onda propagada, las perturbaciones ambientales y las reflexiones producidas por los límites del recinto. La señal eléctrica obtenida después de la transducción incorpora, además, la respuesta en frecuencia y la sensibilidad de cada micrófono, así como las contribuciones del circuito de acondicionamiento y adquisición.

Por tanto, las señales registradas no deben interpretarse directamente como presiones acústicas absolutas sin considerar la calibración de los transductores y de la cadena electrónica. En los apartados siguientes se describen las operaciones que se aplican a las señales adquiridas para obtener observables temporales y espectrales útiles en la caracterización del movimiento de la fuente.

% ==============================================================================
% 1.2. DISCRETIZACIÓN Y ACONDICIONAMIENTO DEL CAMPO ACÚSTICO
% ==============================================================================

\subsection{Discretización y acondicionamiento del campo acústico}

La adquisición transforma las señales eléctricas procedentes de los dos receptores en secuencias discretas que pueden procesarse digitalmente. La cadena comprende la conversión analógico-digital, el promediado de muestras y la formación de bloques para el análisis temporal y espectral. Su organización se resume en la Figura~\ref{fig:discretizacion_flujo}.

\begin{figure}[htbp]
    \centering
    % figures/fig1_discretizacion.tex
\begin{tikzpicture}[
    >=Stealth,
    % Cajas de la columna izquierda (física instrumental)
    block/.style={
        rectangle,
        draw=black!70,
        fill=black!2,
        rounded corners=4pt,
        line width=0.8pt,
        text width=5.8cm,
        align=center,
        inner sep=8pt,
        minimum height=1.35cm,
        font=\small\sffamily
    },
    % Fórmulas matemáticas de la columna derecha
    state/.style={
        font=\small,
        align=center,
        inner sep=6pt
    },
    % Flechas principales
    arrow/.style={
        ->,
        line width=1.05pt,
        draw=black!75,
        shorten >= 2pt,
        shorten <= 2pt
    },
    % Líneas de unión sin punta
    line/.style={
        line width=1.05pt,
        draw=black!75
    },
    % Operadores circulares
    oplabel/.style={
        circle,
        draw=black!85,
        fill=white,
        line width=0.9pt,
        inner sep=3.5pt,
        font=\normalsize\bfseries
    },
    % Cajas terminales de las ramas (Parte 2)
    branchbox/.style={
        rectangle,
        draw=UdeAGreen!85!black,
        fill=UdeAGreen!6,
        rounded corners=4pt,
        line width=0.9pt,
        text width=5.4cm,
        align=center,
        inner sep=7pt,
        font=\small\sffamily
    }
]

    % =========================================================================
    % Nivel 0: Campo Continuo (y = 0.0)
    % =========================================================================
    \node[block] (b0) at (0, 0) {
        \textbf{Campos acústicos continuos}\\
        \footnotesize Presión sonora en receptores duales
    };
    \node[state] (s0) at (7.2, 0) {
        $\mathbf{p}(t) = \begin{pmatrix} p_1(t) \\ p_2(t) \end{pmatrix} \in L^2(\mathbb{R}) \times L^2(\mathbb{R})$
    };

    % Operador 1
    \node[oplabel] (op1) at (7.2, -1.8) {$\hat{\mathcal{P}}_{\Delta t}$};

    % =========================================================================
    % Nivel 1: Muestreo Simultáneo (y = -3.6)
    % =========================================================================
    \node[block] (b1) at (0, -3.6) {
        \textbf{Muestreo simultáneo en FPGA}\\
        \footnotesize Cuantización determinista ($\tau_{\text{skew}} \equiv 0$)
    };
    \node[state] (s1) at (7.2, -3.6) {
        $\mathbf{p}[n] = \begin{pmatrix} p_1[n\Delta t] \pm \Delta p \\ p_2[n\Delta t] \pm \Delta p \end{pmatrix}$
    };

    % Operador 2
    \node[oplabel] (op2) at (7.2, -5.4) {$\hat{\mathcal{C}}_M$};

    % =========================================================================
    % Nivel 2: Diezmado (y = -7.2)
    % =========================================================================
    \node[block] (b2) at (0, -7.2) {
        \textbf{Diezmado digital promediado}\\
        \footnotesize Reducción estocástica de varianza
    };
    \node[state] (s2) at (7.2, -7.2) {
        $\tilde{\mathbf{p}}[k] = \begin{pmatrix} \tilde{p}_1[k\Delta t'] \pm \frac{\Delta p}{\sqrt{M}} \\ \tilde{p}_2[k\Delta t'] \pm \frac{\Delta p}{\sqrt{M}} \end{pmatrix}$
    };

    % Operador 3
    \node[oplabel] (op3) at (7.2, -9.0) {$\hat{\mathcal{W}}_{nc}$};

    % =========================================================================
    % Nivel 3: Enventanado Causal (y = -10.8)
    % =========================================================================
    \node[block] (b3) at (0, -10.8) {
        \textbf{Enventanado temporal causal}\\
        \footnotesize Truncamiento a soporte compacto $\mathbb{R}^N$
    };
    \node[state] (s3) at (7.2, -10.8) {
        $\vec{P} = \begin{pmatrix} \vec{p}_1 \pm \Delta\vec{p}' \\ \vec{p}_2 \pm \Delta\vec{p}' \end{pmatrix} \in \mathbb{R}^N \times \mathbb{R}^N$
    };

    % =========================================================================
    % CONEXIONES VERTICALES INTERNAS (Etapa 1)
    % =========================================================================
    \draw[arrow] (b0) -- (b1);
    \draw[arrow] (b1) -- (b2);
    \draw[arrow] (b2) -- (b3);

    \draw[arrow] (s0) -- (op1);
    \draw[arrow] (op1) -- (s1);

    \draw[arrow] (s1) -- (op2);
    \draw[arrow] (op2) -- (s2);

    \draw[arrow] (s2) -- (op3);
    \draw[arrow] (op3) -- (s3);

    % =========================================================================
    % BIFURCACIÓN ELEGANTE (Convergencia suave y tronco común con etiqueta)
    % =========================================================================
    % 1. Curvas suaves de unión hacia el centro (x = 3.6, y = -12.4)
    \coordinate (merge) at (3.6, -12.4);
    \draw[line] (b3.south) to[out=-90, in=180] (merge);
    \draw[line] (s3.south) to[out=-90, in=0] (merge);

    % 2. Tronco común vertical con etiqueta física del vector que se bifurca
    \coordinate (split) at (3.6, -13.5);
    \draw[line] (merge) -- node[right=5pt, font=\small, align=left] {
        \textbf{Vector de estado $\vec{P}$}\\
        \footnotesize (Acondicionado)
    } (split);

    % 3. Cajas de destino de las ramas (y = -14.8)
    \node[branchbox] (ramaA) at (0.0, -14.8) {
        \textbf{Rama A: Canal temporal}\\
        \footnotesize Banda estrecha en $L^2$ (Sección~1.3)
    };
    \node[branchbox] (ramaB) at (7.2, -14.8) {
        \textbf{Rama B: Espacio recíproco}\\
        \footnotesize Descomposición en $\ell^2$ (Sección~1.4)
    };

    % 4. Flechas ortogonales limpias hacia cada rama
    \draw[arrow] (split) -| (ramaA.north);
    \draw[arrow] (split) -| (ramaB.north);

\end{tikzpicture}
    \vspace{0.4cm}
    \caption{Cadena de adquisición y acondicionamiento de las señales acústicas de los dos receptores. Se indican la conversión analógico-digital, el promediado de muestras y la formación de bloques para el procesamiento temporal y espectral.}
    \label{fig:discretizacion_flujo}
\end{figure}

\subsubsection{Adquisición y muestreo de las señales}

Sean $v_1(t)$ y $v_2(t)$ las señales eléctricas entregadas por los dos canales de recepción. La conversión analógico-digital produce las secuencias:

\begin{equation}
    v_i[n] = v_i(nT_s),
    \qquad i\in\{1,2\},
    \label{eq:operador_muestreo}
\end{equation}

donde $T_s=1/f_s$ es el periodo de muestreo de cada canal. La adquisición descrita para WALO utiliza una frecuencia de muestreo de hasta $\SI{500}{\kilo\hertz}$ por canal antes de la etapa de promediado. Por tanto, el periodo correspondiente es:

\begin{equation}
    T_s = \frac{1}{\SI{500}{\kilo\hertz}}
    = \SI{2}{\micro\second}.
\end{equation}

La cuantización de un convertidor de $12$ bits asigna cada muestra a uno de un conjunto finito de niveles. Si $\Delta V$ representa el paso de cuantización correspondiente al rango de entrada configurado, el error de cuantización puede modelarse, bajo las hipótesis habituales de cuantización uniforme y ausencia de sobrecarga, mediante una variable distribuida en:

\begin{equation}
    -\frac{\Delta V}{2}
    \leq \epsilon_q
    \leq
    \frac{\Delta V}{2}.
\end{equation}

Su varianza teórica es:

\begin{equation}
    \sigma_q^2 = \frac{(\Delta V)^2}{12}.
    \label{eq:varianza_cuantizacion}
\end{equation}

Para un convertidor ideal de $12$ bits, el paso de cuantización depende del rango de tensión de entrada:

\begin{equation}
    \Delta V \approx \frac{V_{\mathrm{rango}}}{2^{12}},
\end{equation}

donde $V_{\mathrm{rango}}$ es el intervalo total de conversión. El valor numérico de la varianza solo puede establecerse una vez especificados el rango efectivo, la referencia y la configuración del convertidor. Además, la expresión~(\ref{eq:varianza_cuantizacion}) describe el modelo ideal de cuantización; no incluye el ruido analógico, la no linealidad del convertidor ni las perturbaciones de la electrónica de entrada \cite{oppenheim1999discrete, lyons2010understanding}.

La adquisición de ambos canales utiliza una temporización común en la lógica programable. Esta arquitectura permite establecer una referencia temporal compartida para las conversiones y evita depender de la secuenciación de lecturas realizada por un procesador convencional. Sin embargo, la existencia de un reloj común no demuestra por sí sola que el desfase temporal físico entre canales sea exactamente nulo. El sesgo residual depende de la temporización efectiva de los convertidores, las rutas de señal y la electrónica asociada, por lo que debe caracterizarse experimentalmente si se requiere una incertidumbre intercanal específica.

La relevancia de esta característica puede estimarse considerando una señal sinusoidal de frecuencia $f_0$. Un desfase temporal $\tau$ entre canales produce una diferencia de fase:

\begin{equation}
    \Delta\phi = 2\pi f_0\tau.
    \label{eq:error_fase_skew}
\end{equation}

Para una portadora de $\SI{2.6}{\kilo\hertz}$, un desfase de $\SI{10}{\micro\second}$ correspondería a:

\begin{equation}
    \Delta\phi
    =
    2\pi
    \left(\SI{2.6}{\kilo\hertz}\right)
    \left(\SI{10}{\micro\second}\right)
    \approx \SI{9.4}{\degree}.
\end{equation}

Este ejemplo muestra por qué la sincronización intercanal es relevante para las mediciones de fase y tiempo de llegada. No constituye, por sí mismo, una medición del desfase residual de WALO.

\subsubsection{Promediado y reducción de la tasa de muestreo}

Después de la adquisición, las muestras pueden agruparse en bloques de $M$ elementos y reemplazarse por su promedio. Para cada canal, la salida se define como:

\begin{equation}
    \tilde{v}_i[k]
    =
    \frac{1}{M}
    \sum_{m=0}^{M-1}
    v_i[Mk-m],
    \qquad i\in\{1,2\}.
    \label{eq:operador_diezmado}
\end{equation}

Esta operación combina un filtrado de respuesta rectangular con una reducción de la tasa de muestreo. El nuevo periodo entre muestras es:

\begin{equation}
    T_s' = MT_s,
\end{equation}

y la frecuencia efectiva por canal resulta:

\begin{equation}
    f_s' = \frac{f_s}{M}.
    \label{eq:frecuencia_diezmada}
\end{equation}

Para el modo nominal con $f_s=\SI{500}{\kilo\hertz}$ por canal y $M=10$, se obtiene:

\begin{equation}
    f_s' = \SI{50}{\kilo\hertz},
    \qquad
    T_s' = \SI{20}{\micro\second}.
\end{equation}

El efecto del promediado sobre el ruido puede analizarse separando la señal en una componente determinista $s_i[n]$ y una perturbación aditiva $\epsilon_i[n]$:

\begin{equation}
    v_i[n] = s_i[n]+\epsilon_i[n].
    \label{eq:descomposicion_ruido}
\end{equation}

Si el ruido tiene media nula, varianza $\sigma^2$ y muestras incorrelacionadas, entonces:

\begin{equation}
    \mathbb{E}[\epsilon_i[n]]=0,
    \qquad
    \mathbb{E}[\epsilon_i[n]\epsilon_i[m]]
    =
    \sigma^2\delta_{nm}.
\end{equation}

En esas condiciones, la varianza del promedio de $M$ muestras es:

\begin{align}
    \operatorname{Var}
    \left(
    \frac{1}{M}\sum_{m=0}^{M-1}\epsilon_i[Mk-m]
    \right)
    &=
    \frac{1}{M^2}
    \sum_{m=0}^{M-1}\sigma^2
    \nonumber\\
    &=
    \frac{\sigma^2}{M}.
    \label{eq:varianza_promediado}
\end{align}

Por tanto, la desviación estándar del ruido se reduce en un factor $\sqrt{M}$ bajo las hipótesis indicadas. Para $M=10$, la reducción teórica de potencia del ruido blanco incorrelacionado es de $\SI{10}{\decibel}$ \cite{lyons2010understanding, mitra2006digital}.

Este resultado no debe interpretarse como una reducción universal de toda perturbación presente en la señal. Si el ruido está correlacionado entre muestras, la varianza del promedio depende de sus correlaciones temporales. Asimismo, el promediado modifica la respuesta en frecuencia de la señal y debe considerarse al evaluar la atenuación de componentes cercanas a la frecuencia de Nyquist de la nueva tasa de muestreo.

En particular, el promedio de $M$ muestras consecutivas no equivale a un filtro antialias ideal. La respuesta de la ventana rectangular presenta lóbulos laterales y no garantiza por sí sola la eliminación de todas las componentes que podrían plegarse durante el submuestreo. La validez de esta etapa depende, por tanto, de la banda de interés y de la evaluación de la respuesta efectiva de la cadena.

\subsubsection{Formación de bloques y procesamiento temporal y espectral}

Para realizar el procesamiento digital, las muestras acondicionadas se agrupan en bloques de longitud $N$. Las longitudes consideradas para el análisis son:

\begin{equation}
    N\in\{512,1024,2048\}.
\end{equation}

Un bloque de cada canal puede representarse como:

\begin{equation}
    \mathbf{v}_i[k_0]
    =
    \begin{pmatrix}
        \tilde{v}_i[k_0] &
        \tilde{v}_i[k_0+1] &
        \cdots &
        \tilde{v}_i[k_0+N-1]
    \end{pmatrix}^{T}.
    \label{eq:bloque_temporal}
\end{equation}

La formación de bloques delimita los intervalos de datos utilizados por los algoritmos posteriores. No implica necesariamente que todas las operaciones deban aplicarse sobre las muestras con la misma ponderación. En particular, las ventanas espectrales se emplean para reducir la fuga espectral en la transformada de Fourier, mientras que la detección de frentes de onda y el análisis de tiempos de llegada requieren conservar la evolución temporal de la señal \cite{oppenheim1999discrete, harris1978use}.

Para el análisis espectral se considera una ventana de Blackman-Harris de cuatro términos, cuya expresión general es:

\begin{equation}
    w[n]
    =
    a_0
    -
    a_1\cos\left(\frac{2\pi n}{N-1}\right)
    +
    a_2\cos\left(\frac{4\pi n}{N-1}\right)
    -
    a_3\cos\left(\frac{6\pi n}{N-1}\right),
    \label{eq:ventana_blackman_harris}
\end{equation}

con coeficientes $a_0,a_1,a_2,a_3$ definidos por la variante de la ventana utilizada. Esta ventana presenta una atenuación elevada de los lóbulos laterales y resulta útil cuando se desea reducir la contribución de componentes espectrales cercanas al tono de interés \cite{harris1978use, oppenheim1999discrete}.

La ventana modifica la amplitud efectiva de las componentes y debe tenerse en cuenta al estimar amplitudes espectrales. Tampoco debe aplicarse indiscriminadamente a todas las operaciones temporales: una ponderación que es conveniente para el análisis de Fourier puede modificar la forma de los frentes de onda y afectar la detección de instantes de llegada.

A partir de los bloques adquiridos se distinguen dos conjuntos de operaciones:
\begin{enumerate}
    \item \textbf{Procesamiento temporal:} filtrado pasabanda, estimación de amplitud y fase instantáneas, cálculo de frecuencia instantánea, detección de tiempos de llegada y acumulación de magnitudes cuadráticas de la señal.

    \item \textbf{Procesamiento espectral:} transformación discreta de Fourier, cálculo de magnitudes y fases espectrales y estimación de la frecuencia de la componente dominante.
\end{enumerate}

Estas operaciones permiten obtener observables complementarios. Su interpretación física y las condiciones bajo las cuales resultan válidos se desarrollan en las secciones siguientes.

% ==============================================================================
% 1.3. RAMA A: REPRESENTACIÓN TEMPORAL Y CANAL PASABANDA
% ==============================================================================

\subsection{Rama A: Representación temporal y canal pasabanda}

El procesamiento temporal tiene como objetivo extraer información sobre la evolución de la señal recibida. En esta rama se consideran tres procedimientos: la estimación de amplitud y frecuencia instantáneas, la acumulación de una magnitud cuadrática en una ventana temporal y la detección de tiempos de llegada.

\subsubsection{Filtrado pasabanda y respuesta de fase}

Las señales adquiridas pueden contener componentes fuera de la banda ocupada por la emisión acústica de interés. Para reducir su contribución se aplica un filtro pasabanda centrado alrededor de la frecuencia portadora $f_0$. La señal filtrada de cada canal se representa como:

\begin{equation}
    y_i[n] = \mathcal{F}_{\mathrm{BP}}\{v_i[n]\},
    \qquad i\in\{1,2\}.
    \label{eq:filtro_fase_cero}
\end{equation}

La respuesta de fase del filtro es relevante cuando se comparan las señales de ambos receptores. Un filtro causal puede introducir un retardo de grupo dependiente de la frecuencia:

\begin{equation}
    \tau_g(\omega)
    =
    -\frac{d\theta(\omega)}{d\omega},
    \label{eq:retardo_grupo}
\end{equation}

donde $\theta(\omega)$ es la fase de la respuesta en frecuencia. Si los canales presentan respuestas diferentes, el filtrado puede modificar la diferencia de fase o de tiempo observada entre ellos \cite{mitra2006digital, oppenheim1999discrete}.

Una alternativa habitual para el procesamiento fuera de línea consiste en aplicar el filtro en dos pasadas, una hacia adelante y otra hacia atrás. Para un filtro con respuesta $H(e^{j\omega})$, la respuesta efectiva ideal de este procedimiento es:

\begin{equation}
    H_{\mathrm{eff}}(e^{j\omega})
    =
    H(e^{j\omega})H(e^{-j\omega})
    =
    \left|H(e^{j\omega})\right|^2,
    \label{eq:respuesta_bidireccional}
\end{equation}

si los coeficientes del filtro son reales. La respuesta resultante es real y no introduce una fase dependiente de la frecuencia en el régimen ideal. Sin embargo, la operación es no causal y puede producir efectos de borde en los extremos de los bloques. Además, modifica la magnitud de la respuesta, que queda elevada al cuadrado respecto de la respuesta de una pasada.

Por estas razones, el filtrado bidireccional no debe interpretarse como una garantía de ausencia de todo error temporal o de fase en la cadena de medición. La coherencia intercanal también depende de que ambos canales utilicen los mismos parámetros de filtrado, de la alineación de los bloques y de la respuesta de los elementos analógicos. La caracterización de estos efectos es necesaria cuando se pretende asignar una incertidumbre a las diferencias de fase o de tiempo.

La arquitectura funcional de esta rama se resume en la Figura~\ref{fig:rama_a_flujo}.

\begin{figure}[htbp]
    \centering
    % figures/fig2_rama_a.tex
\begin{tikzpicture}[
    >=Stealth,
    % Estilo idéntico al bloque principal de la Figura 1
    block/.style={
        rectangle,
        draw=black!70,
        fill=black!2,
        rounded corners=4pt,
        line width=0.8pt,
        text width=6.2cm,
        align=center,
        inner sep=8pt,
        minimum height=1.35cm,
        font=\small\sffamily
    },
    % Estilo idéntico a las cajas terminales de las ramas de la Figura 1
    branchbox/.style={
        rectangle,
        draw=black!70,
        fill=black!4,
        rounded corners=4pt,
        line width=0.9pt,
        text width=4.6cm,
        align=center,
        inner sep=7pt,
        font=\small\sffamily
    },
    % Flechas principales
    arrow/.style={
        ->,
        line width=1.05pt,
        draw=black!75,
        shorten >= 2pt,
        shorten <= 2pt
    },
    % Líneas de unión sin punta
    line/.style={
        line width=1.05pt,
        draw=black!75
    },
    % Operador circular idéntico al de la Figura 1
    oplabel/.style={
        circle,
        draw=black!85,
        fill=white,
        line width=0.9pt,
        inner sep=3.5pt,
        font=\normalsize\bfseries
    }
]

    % =========================================================================
    % Nivel 0: Entrada de la Rama A (y = 0.0)
    % =========================================================================
    \node[block] (inA) at (0, 0) {
        \textbf{Vector de estado $\vec{P}$}\\
        \footnotesize Entrada a la Rama A (Banda estrecha en $L^2$)
    };

    % Operador de Fase Cero (y = -1.8)
    \node[oplabel] (opF) at (0, -1.8) {$\hat{\mathcal{F}}_{\text{ZPh}}$};

    % =========================================================================
    % Nivel 1: Señal Filtrada (y = -3.6)
    % =========================================================================
    \node[block] (filtered) at (0, -3.6) {
        \textbf{Señal filtrada de fase nula}\\
        \footnotesize $\mathbf{y}(t) = \mathbf{y}(t) \pm \Delta\mathbf{y}$ ($\tau_g \equiv 0$)
    };

    % Conexiones superiores
    \draw[arrow] (inA) -- (opF);
    \draw[arrow] (opF) -- (filtered);

    % =========================================================================
    % TRIFURCACIÓN SIMÉTRICA (3 VÍAS)
    % =========================================================================
    % 1. Tronco vertical descendente desde la señal filtrada
    \coordinate (split) at (0, -4.8);
    \draw[line] (filtered.south) -- (split);

    % 2. Barra horizontal de distribución (de -5.4 a +5.4)
    \coordinate (leftbar) at (-5.4, -4.8);
    \coordinate (rightbar) at (5.4, -4.8);
    \draw[line] (leftbar) -- (rightbar);

    % 3. Bajadas verticales hacia las 3 vías (x = -5.4, x = 0.0, x = +5.4)
    \draw[arrow] (-5.4, -4.8) -- (-5.4, -5.6);
    \draw[arrow] (0.0, -4.8) -- (0.0, -5.6);
    \draw[arrow] (5.4, -4.8) -- (5.4, -5.6);

    % =========================================================================
    % Nivel 2: Las 3 Vías de la Rama A (y = -7.2)
    % =========================================================================
    \node[branchbox] (via1) at (-5.4, -7.2) {
        \textbf{Vía 1: Cinemática Doppler}\\
        \footnotesize Transformada de Hilbert $\hat{\mathcal{H}}$\\
        Vector analítico $\mathbf{z}(t) \in \mathbb{C}^N$\\
        Derivada regularizada $\hat{\mathcal{D}}_{\Phi, SG}$
    };

    \node[branchbox] (via2) at (0.0, -7.2) {
        \textbf{Vía 2: Régimen Cuasi-Anecoico}\\
        \footnotesize Enventanado causal $\hat{T}_{\text{gate}}$ ($K=3$)\\
        Demodulación Lock-In $\hat{\mathcal{P}}_{f_0}$\\
        Integral de energía $\hat{\mathcal{E}}$
    };

    \node[branchbox] (via3) at (5.4, -7.2) {
        \textbf{Vía 3: Detección ToA}\\
        \footnotesize Ventana guarda $\hat{\mathcal{G}}_{\text{blank}}$\\
        Discriminador pendiente $\hat{\mathcal{C}}_{zc}$\\
        Interpolador secante $\hat{\mathcal{I}}_{zc}$
    };

\end{tikzpicture}
    \vspace{0.4cm}
    \caption{Procesamiento temporal de las señales recibidas. Se muestran las operaciones de filtrado, estimación de amplitud y frecuencia instantáneas, acumulación de una magnitud cuadrática y detección de tiempos de llegada.}
    \label{fig:rama_a_flujo}
\end{figure}

\subsubsection{Estimación de amplitud y frecuencia instantáneas}

Para una señal real de banda estrecha, la amplitud y la fase pueden estimarse mediante su señal analítica. La transformada de Hilbert de una señal continua $y_i(t)$ se define formalmente como \cite{gabor1946theory, cohen1995time}:

\begin{equation}
    \widetilde{y}_i(t)
    =
    \mathcal{H}\{y_i(t)\}
    =
    \frac{1}{\pi}
    \operatorname{P.V.}
    \int_{-\infty}^{\infty}
    \frac{y_i(\tau)}{t-\tau}\,d\tau,
    \label{eq:hilbert_integral}
\end{equation}

donde $\operatorname{P.V.}$ indica el valor principal de Cauchy. En el procesamiento digital, esta operación se implementa mediante una transformación discreta que depende de la longitud del bloque y de las condiciones de borde. La expresión continua anterior describe la definición matemática, no el algoritmo discreto concreto.

La señal analítica asociada a cada canal es:

\begin{equation}
    z_i(t)
    =
    y_i(t)+j\widetilde{y}_i(t),
    \qquad i\in\{1,2\}.
    \label{eq:vector_analitico}
\end{equation}

A partir de ella se calculan la envolvente y la fase instantáneas:

\begin{equation}
    A_i(t)
    =
    |z_i(t)|
    =
    \sqrt{y_i^2(t)+\widetilde{y}_i^2(t)},
    \label{eq:amplitud_instantanea}
\end{equation}

\begin{equation}
    \Phi_i(t)
    =
    \operatorname{unwrap}
    \left[
    \operatorname{atan2}
    \left(\widetilde{y}_i(t),y_i(t)\right)
    \right].
    \label{eq:fase_desenvuelta}
\end{equation}

La función $\operatorname{unwrap}$ elimina saltos de $2\pi$ entre muestras consecutivas cuando estos corresponden a discontinuidades de representación de la fase. La fase obtenida puede utilizarse para estimar la frecuencia instantánea:

\begin{equation}
    f_i(t)
    =
    \frac{1}{2\pi}
    \frac{d\Phi_i(t)}{dt}.
    \label{eq:frecuencia_derivada}
\end{equation}

En señales discretas, la derivación amplifica las variaciones rápidas de la fase, incluidas aquellas debidas al ruido. Por ello, la derivada directa mediante diferencias finitas puede producir fluctuaciones importantes en la frecuencia estimada.

Una alternativa consiste en ajustar un polinomio local a la fase desenvuelta mediante el procedimiento de Savitzky-Golay \cite{savitzky1964smoothing}. Para un polinomio de grado $p$ ajustado en una ventana simétrica de $2K+1$ muestras, se escribe:

\begin{equation}
    \Phi_i(t_k+\tau)
    \approx
    \sum_{r=0}^{p}c_r\tau^r,
    \qquad
    \tau\in[-KT_s',KT_s'].
    \label{eq:ajuste_savitzky_golay}
\end{equation}

La derivada en el centro de la ventana se obtiene a partir del coeficiente lineal:

\begin{equation}
    f_{i,\mathrm{inst}}(t_k)
    =
    \frac{1}{2\pi}
    \left.
    \frac{d\Phi_i}{dt}
    \right|_{t=t_k}
    =
    \frac{c_1}{2\pi}.
    \label{eq:doppler_regularizado}
\end{equation}

El resultado depende del grado del polinomio, la longitud de la ventana y la calidad de la fase estimada. Una ventana más extensa puede reducir fluctuaciones, pero también suavizar variaciones físicas rápidas. Por tanto, los parámetros del ajuste deben seleccionarse considerando la escala temporal del movimiento que se desea medir.

Las magnitudes $A_i(t)$ y $f_{i,\mathrm{inst}}(t)$ describen, respectivamente, la envolvente y la variación de frecuencia de las señales recibidas. La diferencia entre las frecuencias estimadas en ambos canales puede emplearse posteriormente en un modelo Doppler diferencial. Su interpretación como velocidad requiere especificar la geometría y las hipótesis de propagación correspondientes.

\subsubsection{Acumulación de energía digital en una ventana temporal}

En un recinto cerrado, la señal registrada en cada receptor puede contener tanto la onda directa como contribuciones reflejadas en el suelo, las paredes y el techo. La superposición de estas contribuciones modifica la amplitud resultante y puede producir interferencia constructiva o destructiva \cite{kuttruff2016room}.

Una estrategia para reducir la influencia de reflexiones tardías consiste en limitar el intervalo de análisis a una ventana temporal próxima a la llegada directa. Si la frecuencia nominal de emisión es $f_0$, el periodo correspondiente es:

\begin{equation}
    T_0=\frac{1}{f_0}.
\end{equation}

Para una ventana de $K$ ciclos, su duración es:

\begin{equation}
    T_{\mathrm{gate}}=KT_0=\frac{K}{f_0}.
    \label{eq:tiempo_gate}
\end{equation}

Por ejemplo, con $K=3$ y $f_0=\SI{2.6}{\kilo\hertz}$:

\begin{equation}
    T_{\mathrm{gate}}
    \approx
    \SI{1.15}{\milli\second}.
\end{equation}

La ventana puede reducir la contribución de reflexiones que llegan después de su finalización. No elimina toda la reverberación: una reflexión cuyo tiempo de llegada se encuentre dentro del intervalo seguirá contribuyendo a la señal acumulada. Para justificar el carácter predominantemente directo de la ventana, debe compararse su duración con el tiempo de llegada de la primera reflexión relevante.

Si la diferencia entre la longitud del trayecto reflejado y la del trayecto directo es $\Delta r$, la diferencia de tiempos de llegada correspondiente es:

\begin{equation}
    \Delta t_{\mathrm{refl}}=\frac{\Delta r}{c}.
    \label{eq:tiempo_reflexion}
\end{equation}

La condición para que una reflexión llegue después del intervalo de integración es:

\begin{equation}
    T_{\mathrm{gate}}<\Delta t_{\mathrm{refl}}.
    \label{eq:condicion_gate}
\end{equation}

Esta condición depende de la geometría real del recinto y de la posición de la fuente y los receptores. La distancia absoluta a una pared no equivale necesariamente a la diferencia de trayectos $\Delta r$; por tanto, no basta con especificar una distancia mínima a los límites del laboratorio para demostrar que la ventana está libre de reflexiones.

Dentro de la ventana seleccionada puede estimarse la amplitud de la componente armónica mediante una proyección síncrona sobre la frecuencia de referencia $f_0$. Para $N_{\mathrm{gate}}$ muestras, una expresión discreta del fasor es:

\begin{equation}
    X_i(f_0)
    =
    \frac{2}{N_{\mathrm{gate}}}
    \sum_{m=0}^{N_{\mathrm{gate}}-1}
    y_{i,\mathrm{gate}}[m]
    e^{-j2\pi f_0mT_s'}.
    \label{eq:lock_in_proyeccion}
\end{equation}

La normalización anterior corresponde a una estimación de amplitud para una señal sinusoidal en condiciones apropiadas. Su interpretación cuantitativa depende de la duración de la ventana, del número de ciclos incluidos, de la frecuencia de referencia y de la fase inicial. Si el intervalo no contiene un número entero de ciclos, la proyección puede presentar contribuciones residuales asociadas a la falta de ortogonalidad sobre la ventana finita \cite{meade1983lock}.

Para una componente sinusoidal cuya amplitud se estime correctamente mediante $X_i(f_0)$, la amplitud eficaz puede expresarse como:

\begin{equation}
    A_{\mathrm{RMS},i}
    =
    \frac{|X_i(f_0)|}{\sqrt{2}}.
    \label{eq:amplitud_rms}
\end{equation}

También puede calcularse una suma cuadrática de las muestras dentro de la ventana:

\begin{equation}
    Q_i
    =
    \sum_{m=0}^{N_{\mathrm{gate}}-1}
    \left(
        y_{i,\mathrm{gate}}[m]
        -
        \overline{y}_{i,\mathrm{gate}}
    \right)^2,
    \label{eq:energia_directa}
\end{equation}

donde $\overline{y}_{i,\mathrm{gate}}$ es el valor medio de las muestras consideradas.

La magnitud $Q_i$ es un indicador digital cuadrático de la señal. No constituye por sí sola una energía acústica en unidades físicas ni una densidad de flujo energético. Para obtener magnitudes acústicas absolutas se requiere conocer la sensibilidad del micrófono, la respuesta de la cadena de adquisición, la normalización temporal y las relaciones físicas pertinentes entre presión e intensidad acústica.

\subsubsection{Detección de tiempos de llegada (ToA)}

La medición del tiempo de llegada (*Time of Arrival*, ToA) permite relacionar el instante de recepción de una señal con la distancia recorrida desde la fuente. Para dos receptores, la diferencia entre los tiempos de llegada puede utilizarse en modelos de diferencia temporal (TDOA), habituales en sistemas de localización acústica \cite{brandstein2001microphone}.

En una señal real, la detección requiere establecer un criterio que distinga el frente de onda del ruido y de los transitorios del sistema. Una primera operación consiste en descartar un intervalo inicial asociado a la conmutación del emisor y a los transitorios de adquisición. Se define para ello un tiempo de guarda:

\begin{equation}
    t<t_{\mathrm{blank}}.
    \label{eq:ventana_blank}
\end{equation}

El valor de $t_{\mathrm{blank}}$ debe corresponder a la configuración experimental y al comportamiento medido del sistema. Un valor del orden de $\SI{2}{\milli\second}$ puede utilizarse como parámetro de implementación si se confirma para el montaje concreto; no es una constante general del método.

Una vez superado el intervalo de guarda, puede emplearse un umbral de amplitud para habilitar la detección. Por ejemplo:

\begin{equation}
    V_{\mathrm{th}}
    =
    \max\left(
        V_{\mathrm{noise}},
        \alpha A_{\max}
    \right),
    \label{eq:umbral_toa}
\end{equation}

donde $V_{\mathrm{noise}}$ representa un nivel de referencia asociado al ruido, $A_{\max}$ es una amplitud característica de la señal y $\alpha$ es un factor definido por la implementación. El valor $\alpha=0.20$ debe entenderse como un parámetro de configuración y no como un umbral universal.

Después de habilitar la detección, un criterio posible consiste en identificar un cruce ascendente por cero de la señal filtrada. Si el cruce ocurre entre las muestras $k_0$ y $k_0+1$, se cumple:

\begin{equation}
    y_i[k_0]\leq 0,
    \qquad
    y_i[k_0+1]>0.
    \label{eq:cruce_cero}
\end{equation}

El instante de cruce puede estimarse mediante interpolación lineal entre las dos muestras:

\begin{equation}
    t_{\mathrm{arr},i}
    =
    t_b+
    \left(
        k_0
        -
        \frac{y_i[k_0]}
        {y_i[k_0+1]-y_i[k_0]}
    \right)T_s',
    \label{eq:interpolacion_secante}
\end{equation}

donde $t_b$ es el instante correspondiente al inicio del bloque y $k_0$ es el índice local de la primera muestra del intervalo de cruce. Si el índice utilizado es global respecto de toda la adquisición, se toma $t_b=0$.

Esta interpolación permite obtener un estimador temporal con valores fraccionarios del periodo de muestreo. Sin embargo, no aumenta la frecuencia de muestreo ni garantiza por sí misma una exactitud inferior a $T_s'$. La incertidumbre efectiva depende del ruido, la pendiente del cruce, la distorsión de la forma de onda, la estabilidad de la frecuencia, la respuesta de los canales y el criterio utilizado para seleccionar el cruce.

Con $f_s'=\SI{50}{\kilo\hertz}$, el periodo de muestreo es $\SI{20}{\micro\second}$. La distancia recorrida por el sonido durante ese intervalo es:

\begin{equation}
    cT_s'
    \approx
    \left(\SI{343.2}{\metre\per\second}\right)
    \left(\SI{20}{\micro\second}\right)
    \approx
    \SI{6.86}{\milli\metre}.
\end{equation}

Este valor representa la distancia de propagación correspondiente a un periodo de muestreo, no la incertidumbre completa de la estimación ToA. La interpolación puede mejorar la resolución numérica del instante estimado, pero la precisión metrológica debe establecerse mediante caracterización experimental y un presupuesto de incertidumbre.

% ==============================================================================
% 1.4. RAMA B: REPRESENTACIÓN ESPECTRAL
% ==============================================================================

\subsection{Rama B: Representación espectral y estimación de frecuencia}

La segunda rama de procesamiento analiza las señales en el dominio de la frecuencia. Su objetivo es identificar componentes armónicas y estimar la frecuencia de la señal recibida. Esta información complementa las estimaciones temporales de la sección anterior y puede utilizarse como referencia para los modelos Doppler.

La secuencia de procesamiento se resume en la Figura~\ref{fig:rama_b_flujo}.

\begin{figure}[htbp]
    \centering
    % figures/fig3_rama_b.tex
\begin{tikzpicture}[
    >=Stealth,
    % Estilo idéntico al bloque principal de Figuras 1 y 2
    block/.style={
        rectangle,
        draw=black!70,
        fill=black!2,
        rounded corners=4pt,
        line width=0.8pt,
        text width=5.8cm,
        align=center,
        inner sep=8pt,
        minimum height=1.35cm,
        font=\small\sffamily
    },
    % Fórmulas matemáticas de la columna derecha
    state/.style={
        font=\small,
        align=center,
        inner sep=6pt
    },
    % Flechas principales
    arrow/.style={
        ->,
        line width=1.05pt,
        draw=black!75,
        shorten >= 2pt,
        shorten <= 2pt
    },
    % Operador circular idéntico a Figuras 1 y 2
    oplabel/.style={
        circle,
        draw=black!85,
        fill=white,
        line width=0.9pt,
        inner sep=3.5pt,
        font=\normalsize\bfseries
    }
]

    % =========================================================================
    % Nivel 0: Entrada de la Rama B (y = 0.0)
    % =========================================================================
    \node[block] (inB) at (0, 0) {
        \textbf{Vector de estado $\vec{P}$}\\
        \footnotesize Entrada a la Rama B (Espacio recíproco en $\ell^2$)
    };
    \node[state] (s0) at (7.2, 0) {
        $\vec{P} = \begin{pmatrix} \vec{p}_1 \\ \vec{p}_2 \end{pmatrix} \in \mathbb{R}^N \times \mathbb{R}^N$
    };

    % Operador 1
    \node[oplabel] (op1) at (7.2, -1.8) {$\hat{\mathcal{F}}_{\text{FFT}}$};

    % =========================================================================
    % Nivel 1: Espectro Cartesiano (y = -3.6)
    % =========================================================================
    \node[block] (b1) at (0, -3.6) {
        \textbf{Transformada discreta de Fourier}\\
        \footnotesize Vectores espectrales cartesianos $\{X_R, X_I\}$
    };
    \node[state] (s1) at (7.2, -3.6) {
        $\mathbf{Z}[k] = X_R[k] + j X_I[k] \in \mathbb{C}^{N/2}$
    };

    % Operador 2
    \node[oplabel] (op2) at (7.2, -5.4) {$\hat{\mathcal{R}}_C$};

    % =========================================================================
    % Nivel 2: Espectro Polar / CORDIC (y = -7.2)
    % =========================================================================
    \node[block] (b2) at (0, -7.2) {
        \textbf{Rotaciones polares CORDIC}\\
        \footnotesize Transformación afín de Givens
    };
    \node[state] (s2) at (7.2, -7.2) {
        $\left( \norm{\mathbf{Z}[k]}, \, \Phi[k] \right)$
    };

    % Operador 3
    \node[oplabel] (op3) at (7.2, -9.0) {$\hat{\mathcal{M}}_{\text{sinc}}$};

    % =========================================================================
    % Nivel 3: Superresolución Portadora (y = -10.8)
    % =========================================================================
    \node[block] (b3) at (0, -10.8) {
        \textbf{Estimación analítica sub-Hertz}\\
        \footnotesize Inversión no lineal del lóbulo sinc
    };
    \node[state] (s3) at (7.2, -10.8) {
        $\left( f_0 \pm \Delta f_0, \, \norm{\mathbf{Z}(f_0)}, \, \Phi(f_0) \right)$
    };

    % =========================================================================
    % CONEXIONES VERTICALES
    % =========================================================================
    \draw[arrow] (inB) -- (b1);
    \draw[arrow] (b1) -- (b2);
    \draw[arrow] (b2) -- (b3);

    \draw[arrow] (s0) -- (op1);
    \draw[arrow] (op1) -- (s1);

    \draw[arrow] (s1) -- (op2);
    \draw[arrow] (op2) -- (s2);

    \draw[arrow] (s2) -- (op3);
    \draw[arrow] (op3) -- (s3);

\end{tikzpicture}
    \vspace{0.4cm}
    \caption{Procesamiento espectral de las señales adquiridas. La transformada discreta de Fourier proporciona las componentes espectrales complejas, a partir de las cuales se calculan magnitudes y fases y se identifica la frecuencia dominante.}
    \label{fig:rama_b_flujo}
\end{figure}

\subsubsection{Transformada discreta de Fourier y cálculo de magnitud y fase}

Para un bloque de $N$ muestras, la transformada discreta de Fourier (DFT) se define como \cite{cooley1965algorithm, oppenheim1999discrete}:

\begin{equation}
    Z_i[k]
    =
    \sum_{m=0}^{N-1}
    v_i[m]e^{-j2\pi km/N},
    \qquad
    k=0,\ldots,N-1.
    \label{eq:transformada_fourier}
\end{equation}

La transformada rápida de Fourier (FFT) permite calcular esta transformación con menor costo computacional. Su salida es compleja y puede escribirse como:

\begin{equation}
    Z_i[k]=X_{R,i}[k]+jX_{I,i}[k].
\end{equation}

La magnitud y la fase correspondientes son:

\begin{equation}
    |Z_i[k]|
    =
    \sqrt{X_{R,i}^2[k]+X_{I,i}^2[k]},
    \label{eq:modulo_espectral}
\end{equation}

\begin{equation}
    \Phi_i[k]
    =
    \operatorname{atan2}
    \left(
        X_{I,i}[k],X_{R,i}[k]
    \right).
    \label{eq:fase_espectral}
\end{equation}

La implementación puede utilizar el algoritmo CORDIC para calcular magnitud y fase a partir de las componentes cartesianas \cite{volder1959cordic}. Este algoritmo realiza las operaciones mediante rotaciones sucesivas y resulta apropiado para implementaciones digitales en lógica programable. La precisión obtenida depende de la representación numérica, el número de iteraciones y la configuración del bloque de hardware.

La fase espectral debe interpretarse teniendo en cuenta el origen temporal del bloque, la ventana aplicada y la respuesta de los canales. Por tanto, una diferencia de fase entre dos bines o dos canales no puede atribuirse exclusivamente a la propagación acústica sin considerar esos factores instrumentales.

\subsubsection{Resolución espectral e interpolación de la frecuencia dominante}

Para un bloque de $N$ muestras adquirido a una frecuencia de muestreo $f_s'$, la separación entre bines de la DFT es:

\begin{equation}
    \Delta f=\frac{f_s'}{N}.
    \label{eq:resolucion_bin}
\end{equation}

Con una frecuencia de muestreo de $\SI{50}{\kilo\hertz}$ y un bloque de $N=2048$ muestras:

\begin{equation}
    \Delta f
    =
    \frac{\SI{50}{\kilo\hertz}}{2048}
    \approx
    \SI{24.41}{\hertz}.
\end{equation}

La separación entre bines determina el muestreo del espectro producido por la DFT. No representa, en todos los casos, el límite absoluto de precisión con el que puede estimarse la frecuencia de una señal. Bajo condiciones apropiadas de relación señal-ruido, estabilidad de la fuente y modelo espectral, es posible estimar la posición de un máximo entre dos bines mediante interpolación \cite{grandke1983interpolation, jacobsen2007fast}.

Sea $k_0$ el índice del bin de mayor magnitud en la región de interés. Una estimación interpolada puede construirse a partir de las magnitudes de los bines vecinos. En el caso de una interpolación parabólica local sobre magnitudes, por ejemplo, el desplazamiento fraccionario se estima mediante:

\begin{equation}
    \delta
    =
    \frac{1}{2}
    \frac{
        |Z[k_0-1]|-|Z[k_0+1]|
    }{
        |Z[k_0-1]|-2|Z[k_0]|+|Z[k_0+1]|
    }.
    \label{eq:interpolacion_parabolica}
\end{equation}

La frecuencia estimada es entonces:

\begin{equation}
    \widehat{f}_0
    =
    (k_0+\delta)\frac{f_s'}{N}.
    \label{eq:frecuencia_sub_hertz}
\end{equation}

La expresión anterior es una interpolación local y debe utilizarse dentro de las condiciones en las que la aproximación parabólica del máximo espectral sea adecuada. Su sesgo depende de la ventana, la posición fraccionaria del tono, la presencia de componentes cercanas y la relación señal-ruido.

También existen estimadores basados en modelos de la forma del lóbulo principal de la ventana, incluidos métodos de interpolación sinc \cite{grandke1983interpolation, jacobsen2007fast}. Sin embargo, la forma del lóbulo depende de la ventana utilizada. Por ello, no debe aplicarse una fórmula de interpolación sinc sin comprobar que corresponde a la ventana y a la definición de magnitud espectral empleadas en la implementación.

La interpolación se realiza sobre magnitudes lineales o sobre una transformación expresamente considerada por el estimador. No es correcto trasladar directamente una fórmula deducida para magnitudes lineales a datos en decibelios sin analizar el efecto de la transformación logarítmica.

La estimación entre bines permite obtener una frecuencia con más cifras que la separación nominal $\Delta f$, pero esas cifras adicionales no equivalen automáticamente a una incertidumbre menor. La incertidumbre de $\widehat{f}_0$ debe determinarse mediante análisis del estimador y ensayos con señales de frecuencia conocida. En consecuencia, una incertidumbre específica —por ejemplo, del orden de centésimas de hercio— solo puede afirmarse si está respaldada por una caracterización experimental bajo condiciones comparables a las de uso.

% ==============================================================================
% 1.5. SÍNTESIS DE OBSERVABLES Y ACOPLAMIENTO CINEMÁTICO
% ==============================================================================

\subsection{Síntesis de observables y acoplamiento cinemático}

Las etapas anteriores producen magnitudes temporales y espectrales asociadas a cada receptor. Entre ellas se encuentran la amplitud instantánea, la frecuencia instantánea, la amplitud eficaz, la suma cuadrática de la señal, el tiempo de llegada y la frecuencia dominante estimada.

Para organizar estas magnitudes se define el conjunto de observables:

\begin{equation}
    \mathcal{O}(t)
    =
    \left\{
        A_i(t),
        f_{i,\mathrm{inst}}(t),
        A_{\mathrm{RMS},i},
        Q_i,
        t_{\mathrm{arr},i},
        \widehat{f}_{0,i},
        \Phi_i
    \right\}_{i=1,2}.
    \label{eq:vector_observables_matriz}
\end{equation}

No todos los elementos tienen necesariamente la misma dependencia temporal. Las amplitudes y frecuencias instantáneas pueden evaluarse muestra a muestra o por bloques, mientras que las magnitudes RMS, las sumas cuadráticas y las frecuencias dominantes suelen asociarse al intervalo de análisis utilizado. La fase también depende de la referencia temporal y del procedimiento de estimación.

Esta distinción es importante al combinar observables: cada magnitud debe conservar su unidad, su intervalo de estimación y las condiciones bajo las cuales se obtuvo.

\subsubsection{Propagación de incertidumbre y covarianza instrumental}

La incertidumbre de los observables depende tanto de las señales adquiridas como de los procedimientos de procesamiento. La Guía para la Expresión de la Incertidumbre de Medida (GUM) establece un marco para propagar incertidumbres a través de modelos de medición \cite{jcg2008evaluation, taylor1997introduction}.

Si un conjunto de magnitudes de entrada se representa mediante el vector $\mathbf{x}$ y los observables se calculan mediante una función diferenciable:

\begin{equation}
    \mathbf{O}=\mathbf{f}(\mathbf{x}),
\end{equation}

la propagación de incertidumbre en la aproximación lineal de primer orden se expresa mediante:

\begin{equation}
    \mathbf{C}_{O}
    =
    \mathbf{J}
    \mathbf{C}_{x}
    \mathbf{J}^{T},
    \label{eq:propagacion_gum}
\end{equation}

donde $\mathbf{C}_{x}$ es la matriz de covarianza de las magnitudes de entrada y $\mathbf{J}$ es la matriz jacobiana del modelo:

\begin{equation}
    J_{km}
    =
    \frac{\partial f_k}{\partial x_m}.
\end{equation}

Para un observable individual $O_k$, su incertidumbre estándar combinada corresponde a la raíz cuadrada del elemento diagonal asociado:

\begin{equation}
    u(O_k)=\sqrt{(\mathbf{C}_{O})_{kk}}.
\end{equation}

La matriz de covarianza permite considerar tanto las incertidumbres individuales como las correlaciones entre muestras y entre canales. Estas correlaciones pueden aparecer por el ruido compartido, el filtrado, el promediado, la ventana utilizada y los procedimientos de estimación.

No puede suponerse que los términos de covarianza se anulan únicamente porque dos operaciones utilicen bases ortogonales. La ortogonalidad matemática de ciertas funciones no demuestra la independencia estadística de los errores instrumentales ni la ausencia de correlación entre los observables calculados. Si no se dispone de una caracterización de las covarianzas, estas deben identificarse como una contribución pendiente de evaluación y no como términos despreciables por construcción.

Del mismo modo, la resolución nominal del convertidor, el periodo de muestreo o la frecuencia del reloj de la FPGA no determinan por sí solos la incertidumbre de los observables. La incertidumbre efectiva debe incluir, según corresponda, las contribuciones de los transductores, la electrónica analógica, la adquisición, la temporización, el procesamiento y la repetibilidad experimental.

\subsubsection{Relación entre observables y modelos físicos}

Los observables obtenidos en este capítulo constituyen las magnitudes de entrada para los modelos geométricos y cinemáticos desarrollados en el capítulo siguiente. Su relación conceptual se resume en la Tabla~\ref{tab:acoplamiento_observables}.

\begin{table}[htbp]
    \centering
    \small
    \renewcommand{\arraystretch}{1.2}
    \begin{tabularx}{\textwidth}{
        >{\raggedright\arraybackslash}p{0.23\textwidth}
        X
        >{\raggedright\arraybackslash}p{0.25\textwidth}
    }
        \toprule
        \textbf{Observable} &
        \textbf{Modelo físico asociado} &
        \textbf{Magnitud estimable} \\
        \midrule
        Frecuencias instantáneas
        $f_{1,\mathrm{inst}},f_{2,\mathrm{inst}}$
        &
        Efecto Doppler y diferencia de frecuencia entre receptores
        &
        Velocidad o componente radial de velocidad
        \\
        \addlinespace
        Tiempos de llegada
        $t_{\mathrm{arr},1},t_{\mathrm{arr},2}$
        &
        Propagación temporal y diferencia de tiempos de llegada
        &
        Distancia o posición, según la geometría
        \\
        \addlinespace
        Fases de las señales
        $\Phi_1,\Phi_2$
        &
        Diferencia de fase e interferometría acústica
        &
        Dirección de llegada, bajo las condiciones de no ambigüedad
        \\
        \addlinespace
        Amplitudes eficaces
        $A_{\mathrm{RMS},1},A_{\mathrm{RMS},2}$
        &
        Atenuación de la onda y calibración de amplitud
        &
        Distancia, si se conoce el modelo de propagación
        \\
        \addlinespace
        Sumas cuadráticas
        $Q_1,Q_2$
        &
        Comparación de niveles relativos de señal
        &
        Indicador de variación de amplitud o energía digital
        \\
        \addlinespace
        Frecuencias dominantes
        $\widehat{f}_{0,1},\widehat{f}_{0,2}$
        &
        Análisis espectral y estimación de desplazamiento de frecuencia
        &
        Frecuencia de emisión o frecuencia recibida
        \\
        \bottomrule
    \end{tabularx}
    \vspace{0.2cm}
    \caption{Relación entre los observables extraídos de las señales acústicas y los modelos físicos que pueden emplearse para estimar magnitudes geométricas y cinemáticas. La correspondencia depende de las hipótesis de propagación, la geometría del arreglo y la calibración instrumental.}
    \label{tab:acoplamiento_observables}
\end{table}

La relación entre una magnitud medida y una magnitud física estimada no es necesariamente unívoca. Por ejemplo, la amplitud recibida depende de la distancia, pero también de la potencia acústica emitida, la directividad de la fuente, la orientación de los receptores, la respuesta de los micrófonos y las condiciones del recinto. De manera análoga, una diferencia de fase puede relacionarse con la dirección de llegada, pero su interpretación requiere considerar la longitud de onda, la separación entre receptores y las posibles ambigüedades espaciales.

En consecuencia, los observables de este capítulo no determinan por sí solos la posición o la velocidad de la fuente. Proporcionan las magnitudes experimentales que, combinadas con modelos de propagación y restricciones geométricas, permiten formular los problemas de inversión tratados en el capítulo siguiente.
